% ============================================================
% RAPPORT DE PROJET - ENSA Guide
% Plateforme Web d'Orientation Académique
% ============================================================

\documentclass[a4paper,12pt]{article}

% ---- Encodage et langue ----
\usepackage[utf8]{inputenc}
\usepackage[T1]{fontenc}
\usepackage{lmodern}
\usepackage[french]{babel}
\usepackage{csquotes}
\usepackage{placeins}

% ---- Mathématiques ----
\usepackage{amsmath, amssymb}
%--------------------- Variables rapport -----------------
\newcommand{\TitreRapport}{}
\newcommand{\SujetRapport}{}
\newcommand{\ModuleRapport}{}
\newcommand{\Encadrant}{}
\newcommand{\Etudiant}{}
\newcommand{\Jury}{}
\newcommand{\SousTitreRapport}{}

% ---- Formatage des titres (Chapitres) ----
\usepackage{titlesec}
\titleformat{\section}[block]
  {\normalfont\LARGE\bfseries}
  {Chapitre \thesection~:}
  {0.4em}
  {}

\usepackage{titletoc}
\titlecontents{section}[0em]
  {\vspace{0.3cm}\bfseries}
  {Chapitre \thecontentslabel~: }
  {}
  {\titlerule*[0.7pc]{.}\contentspage}

% ---- Graphiques et tableaux ----
\usepackage{graphicx}
\usepackage{float}
\usepackage{array}
\usepackage{booktabs}
\usepackage[a4paper, margin=3.25cm]{geometry}
\emergencystretch=3em
\hbadness=10000
\hfuzz=32pt
\vfuzz=17pt

% ---- Listes ----
\usepackage{enumitem}

% ---- TikZ et diagrammes ----
\usepackage{tikz}
\usepackage{pgfgantt}
\usetikzlibrary{arrows.meta, positioning}

% ---- Hyperliens ----
\usepackage{hyperref}
\hypersetup{
    colorlinks=true,
    linkcolor=black,
    urlcolor=black,
    citecolor=black
}
\usepackage{xurl}
% ---- Code source ----
\usepackage{listings}
\usepackage{xcolor}
\lstdefinestyle{filetree}{
    basicstyle=\ttfamily\small,
    backgroundcolor=\color{brownlight},
    frame=single,
    framerule=1pt,
    rulecolor=\color{brownmain},
    xleftmargin=12pt,
    xrightmargin=6pt,
    framexleftmargin=8pt,
    framexrightmargin=4pt,
    framextopmargin=6pt,
    framexbottommargin=6pt,
    numbers=none,
    breaklines=false,
    keepspaces=true,
    showstringspaces=false,
    commentstyle=\color{brownmid}\itshape,
    morecomment=[l]{\#},
}

\definecolor{codegray}{gray}{0.95}
\definecolor{codeblue}{RGB}{0,0,200}
\definecolor{codegreen}{RGB}{0,120,30}
\definecolor{codered}{RGB}{180,0,0}

\lstset{
    language=Python,
    backgroundcolor=\color{codegray},
    basicstyle=\ttfamily\small,
    keywordstyle=\color{codeblue},
    commentstyle=\color{codegreen!80!black},
    stringstyle=\color{codered},
    numbers=left,
    numberstyle=\tiny\color{gray},
    frame=single,
    breaklines=true,
    showstringspaces=false,
    tabsize=4
}

\lstdefinelanguage{SQL}{
    keywords={SELECT, INSERT, UPDATE, DELETE, FROM, WHERE, CREATE, TABLE,
               INTO, VALUES, DROP, IF, EXISTS, PRIMARY, KEY, NOT, NULL,
               UNIQUE, INTEGER, TEXT, AUTOINCREMENT, FOREIGN, REFERENCES,
               DEFAULT, DATETIME},
    keywordstyle=\color{codeblue},
    sensitive=false,
    comment=[l]{--},
    commentstyle=\color{codegreen},
    stringstyle=\color{codered},
    morestring=[b]',
    morestring=[b]" % chktex 18
}
% ---- Couleurs du thème brun ----
\definecolor{brownmain}{RGB}{107,58,42}
\definecolor{brownlight}{HTML}{FDF6F0}
\definecolor{brownmid}{HTML}{A0785A}

% ---- Graphiques et tableaux ----
\usepackage{graphicx}
\usepackage{array}
\usepackage{booktabs}
\usepackage{tabularx}
\usepackage{colortbl}
\usepackage{caption}
\newcolumntype{L}[1]{>{\raggedright\arraybackslash}p{#1}}

% ---- Boîtes stylisées ----
\usepackage{tcolorbox}
\tcbuselibrary{skins}

\newtcolorbox{infobox}{
  enhanced,
  colback=brownlight,
  colframe=brownmain,
  leftrule=4pt,
  toprule=0.5pt,
  bottomrule=0.5pt,
  rightrule=0.5pt,
  arc=2pt,
  left=8pt, right=8pt, top=6pt, bottom=6pt,
}

% ---- Graphiques pgf ----
\usepackage{pgfplots}
\pgfplotsset{compat=1.18}

% ---- En-tête et pied de page ----
\usepackage{fancyhdr}
\pagestyle{fancy}
\fancyhf{}
\fancyfoot[C]{\thepage}
\renewcommand{\headrulewidth}{0pt}

%--------------------- Page de garde --------------------
\newcommand{\fairepagedegarde}{
\begin{titlepage}
    \centering
    % En-tête avec les deux logos
    \begin{minipage}{0.45\textwidth}
        \begin{flushleft}
            \includegraphics[height=4.5cm]{images/ENSA.png}
        \end{flushleft}
    \end{minipage}
    \hfill
    \begin{minipage}{0.45\textwidth}
        \begin{flushright}
            \includegraphics[height=4.5cm]{images/Caadi.png}
        \end{flushright}
    \end{minipage}

    \vspace{0.4cm}
    {\scshape\Large Université Cadi Ayyad \\[0.2cm] École Nationale des Sciences Appliquées de Safi \\[0.2cm] Département Informatique, Réseaux et Télécommunications \par}

    \vspace{0.9cm}
    {\scshape\Large \SujetRapport \par}

    \vspace{0.9cm}
    {\rule{0.9\linewidth}{1.1pt}\par}
    \vspace{0.45cm}
    {\Huge\bfseries \TitreRapport \par}
    \vspace{0.35cm}
    {\fontsize{14.4}{18}\selectfont\itshape \SousTitreRapport\/ \par}
    \vspace{0.45cm}
    {\rule{0.9\linewidth}{1.1pt}\par}

    \vspace{0.8cm}
    \begin{minipage}[t]{0.4\textwidth}
        \begin{flushleft} \Large
            \textit{\textbf{Réalisé par:}}\par
            \vspace{0.3cm}
            \large \Etudiant
        \end{flushleft}
    \end{minipage}
    \hfill
    \begin{minipage}[t]{0.5\textwidth}
        \begin{flushright} \Large
            \textit{\textbf{Encadré à l'ENSA par:}}\par
            \vspace{0.3cm}
            \large \Encadrant
        \end{flushright}
    \end{minipage}

    \vspace{0.8cm}
    {\large Soutenu le \textbf{20/05/2026} devant le jury composé de: \par}
    \vspace{0.4cm}
    \begin{center}
    \large
    \begin{tabular}{ll}
        \Jury
    \end{tabular}
    \end{center}

    \vfill
    {\Large Année Universitaire: \textbf{2025--2026} \par}
\end{titlepage}
}

%--------------------- Début document ------------------
\title{Rapport de Projet Académique}
\date{}

\begin{document}
\hypersetup{pageanchor=false}

% Different margins for cover page
\newgeometry{
    top=2cm,
    bottom=3.5cm,
    left=2.5cm,
    right=2.5cm
}
%----------- Informations du rapport ---------
\renewcommand{\TitreRapport}{Plateforme Web\\d'Orientation Académique}
\renewcommand{\SousTitreRapport}{EnsaGuide -- Système intelligent d'aide à l'orientation}
\renewcommand{\SujetRapport}{Rapport du mini projet}
\renewcommand{\Encadrant}{M. OUAFIK Youssef}
\renewcommand{\Etudiant}{EL KIHEL Riyad \\ EL MESSAOUDI Othmane \\ HALOUBI Ahmed}
\renewcommand{\Jury}{\textbf{Mme. MESTOURI Hind} & Professeur à l'ENSA de Safi \\ \textbf{M. OUAFIK Youssef} & Professeur à l'ENSA de Safi \\}

%----------- Page de garde -----------------
\fairepagedegarde

% Return to normal margins
\restoregeometry


\pagenumbering{roman}

% =======================
% Dédicaces
% =======================
\newpage
\addcontentsline{toc}{section}{Dédicaces}
\section*{\Huge{Dédicaces}}
\vspace{2cm}
\begin{center}
\textit{
À nos chers parents, pour leurs sacrifices, leur soutien inconditionnel et leurs prières continues.\\[0.5cm]
À nos familles, nos amis et à tous ceux qui ont cru en nous et nous ont encouragés tout au long de notre parcours d'études.\\[0.5cm]
À nos professeurs pour leur dévouement.\\[0.5cm]
À tous nos camarades de l'École Nationale des Sciences Appliquées de Safi.\\[0.5cm]
Nous dédions ce travail.
}
\end{center}

% =======================
% Remerciements
% =======================
\newpage
\addcontentsline{toc}{section}{Remerciements}
\section*{\Huge{Remerciements}}
\vspace{1cm}
{\large \itshape
Nous tenons tout d'abord à exprimer notre sincère gratitude au Professeur
{\normalsize\textbf{OUAFIK Youssef}} pour son encadrement, ses précieux conseils et sa
disponibilité tout au long de ce projet. Son expertise et son encouragement
ont été essentiels pour mener à bien notre travail.

\vspace{0.5cm}

Nous souhaitons également adresser nos chaleureux remerciements aux membres
du jury, {\normalsize\textbf{Mme. MESTOURI Hind}} et {\normalsize\textbf{M. OUAFIK Youssef}}, pour le
temps qu'ils ont consacré à l'évaluation de ce travail et pour leurs remarques
constructives.

\vspace{0.5cm}

Nous souhaitons aussi adresser nos remerciements au personnel enseignant et
administratif de l'École Nationale des Sciences Appliquées de Safi, dont
l'enseignement de qualité nous a permis d'acquérir les connaissances nécessaires
pour réaliser ce projet avec succès.

\vspace{0.5cm}

Nos remerciements vont également à l'endroit de nos amis et collègues, dont
le soutien moral et intellectuel a été inestimable. Leurs discussions stimulantes
et leurs encouragements ont contribué à enrichir notre réflexion et à surmonter
les obstacles rencontrés.

\vspace{0.5cm}

Enfin, nous, {\normalsize\textbf{EL KIHEL Riyad}}, {\normalsize\textbf{EL MESSAOUDI Othmane}} et
{\normalsize\textbf{HALOUBI Ahmed}}, exprimons notre profonde gratitude à nos parents et
à nos proches pour leur soutien indéfectible, leurs conseils avisés et leur
compréhension tout au long de cette période.

\vspace{0.5cm}

Ce projet n'aurait pas été possible sans l'appui et la collaboration de toutes
ces personnes, et nous leur en sommes infiniment reconnaissants.
}

% =======================
% Résumé
% =======================
\newpage
\addcontentsline{toc}{section}{Résumé}
\section*{\Huge{Résumé}}
\vspace{1cm}
{\large
À l'issue du cycle préparatoire de l'École Nationale des Sciences Appliquées de Safi, les étudiants sont confrontés au choix crucial de leur filière d'ingénierie parmi les six options proposées. Ce projet a pour \textbf{objectif} de concevoir et de développer \enquote{ENSA Guide}, une plateforme web interactive d'aide à l'orientation. La \textbf{méthodologie} adoptée repose sur le développement d'une architecture MVC en Python (Flask) couplée à une base de données SQLite. L'outil intègre un questionnaire psychométrique dimensionnel avec calcul de compatibilité (Z-score et Softmax) ainsi qu'un assistant IA conversationnel via l'API Mistral. Les principales \textbf{conclusions} démontrent que la plateforme permet une évaluation précise, rapide et objective des profils étudiants, réduisant ainsi l'incertitude et optimisant leur processus de choix de filière.

\vspace{0.5cm}
\textbf{Mots-clés :} Orientation académique, ENSA Safi, Développement Web, Algorithme de recommandation, Flask, Intelligence Artificielle.
}

% =======================
% Abstract
% =======================
\newpage
\addcontentsline{toc}{section}{Abstract}
\section*{\Huge{Abstract}}
\vspace{1cm}
{\large
At the end of the preparatory cycle at the National School of Applied Sciences (ENSA) of Safi, students face the crucial choice of selecting their engineering specialty from six available options. The \textbf{objective} of this project is to design and develop \enquote{ENSA Guide,} an interactive web platform for academic orientation. The adopted \textbf{methodology} is based on the development of an MVC architecture in Python (Flask) coupled with an SQLite database. The tool integrates a dimensional psychometric questionnaire with a compatibility calculation algorithm (Z-score and Softmax) as well as a conversational AI assistant via the Mistral API. The main \textbf{conclusions} demonstrate that the platform allows for accurate, fast, and objective evaluations of student profiles, thereby reducing uncertainty and optimizing their specialty selection process.

\vspace{0.5cm}
\textbf{Keywords :} Academic orientation, ENSA Safi, Web Development, Recommendation Algorithm, Flask, Artificial Intelligence.
}

% =======================
% Liste des abréviations
% =======================
\newpage
\addcontentsline{toc}{section}{Liste des abréviations}
\section*{\Huge{Liste des abréviations}}
\vspace{1cm}
\begin{itemize}[label=\textbullet, font=\large\bfseries, itemsep=0.5cm]
    \item \textbf{ENSA} : École Nationale des Sciences Appliquées
    \item \textbf{GIIA} : Génie Informatique et Intelligence Artificielle
    \item \textbf{GTR} : Génie des Télécommunications et Réseaux
    \item \textbf{GINDUS} : Génie Industriel
    \item \textbf{GMSI} : Génie Mécatronique et Systèmes Intelligents
    \item \textbf{GPMA} : Génie des Procédés et Matériaux Avancés
    \item \textbf{GATE} : Génie Aéronautique et Technologies de l'Espace
    \item \textbf{MVC} : Modèle-Vue-Contrôleur
    \item \textbf{IA} : Intelligence Artificielle
    \item \textbf{API} : Application Programming Interface
    \item \textbf{JSON} : JavaScript Object Notation
\end{itemize}

% =======================
% Listes et Table des matières
% =======================
\newpage
\addcontentsline{toc}{section}{Liste des figures}
\renewcommand{\listfigurename}{Liste des figures}
\listoffigures

\newpage
\addcontentsline{toc}{section}{Liste des tableaux}
\renewcommand{\listtablename}{Liste des tableaux}
\listoftables

\newpage
\addcontentsline{toc}{section}{Table des matières}
\renewcommand{\contentsname}{Table des matières}
\tableofcontents
\newpage
\clearpage
\hypersetup{pageanchor=true}
\pagenumbering{arabic}
% ============================================================
\clearpage
\section{Introduction générale}
% ============================================================

À l'issue du cycle préparatoire de l'École Nationale des Sciences Appliquées (ENSA) de Safi, les étudiants font face à une décision académique décisive : le choix de leur spécialité d'ingénierie parmi les six filières proposées par l'établissement (GIIA, GTR, GINDUS, GMSI, GPMA et GATE). Cette étape conditionne de manière irréversible leur trajectoire professionnelle et exige une réflexion approfondie.

Cependant, l'absence d'une plateforme d'information centralisée et structurée complexifie grandement cette démarche. Les informations relatives aux maquettes pédagogiques, aux compétences visées et aux débouchés sont souvent fragmentées, ce qui engendre une \textbf{asymétrie d'information structurelle}. L'étudiant se retrouve démuni face à la difficulté d'évaluer l'adéquation entre ces formations et son profil cognitif.

Pour pallier cette carence, le projet \textbf{ENSA Guide} a été initié avec pour ambition de concevoir une plateforme web interactive d'aide à l'orientation. L'objectif fondamental est de fournir un accompagnement technologique sur mesure via un test psychométrique dimensionnel rigoureux et un assistant virtuel conversationnel, garantissant ainsi une prise de décision éclairée, objective et personnalisée.

Afin d'exposer de manière méthodique l'architecture et la réalisation de cette plateforme, ce rapport s'articule autour des axes suivants :
\begin{itemize}
    \item \textbf{Chapitre 2 : Méthodologie et Architecture} -- Présente l'approche incrémentale du projet, l'algorithmique du scoring de compatibilité et l'organisation du modèle MVC.
    \item \textbf{Chapitre 3 : Outils et technologies} -- Expose la pile technologique déployée (backend, base de données, frontend) et l'intégration de l'intelligence artificielle.
    \item \textbf{Chapitre 4 : Implémentation} -- Détaille la conception du schéma relationnel, l'authentification et l'ingénierie du module d'assistance Mistral AI.
    \item \textbf{Chapitre 5 : Interface} -- Illustre le design, l'expérience de navigation de l'étudiant et l'ergonomie du tableau de bord d'administration.
\end{itemize}
% ============================================================
\clearpage
\section{Méthodologie et Architecture}
% ============================================================
\subsection{Approche de développement}

Le projet a suivi une approche de développement incrémentale : les fonctionnalités ont été construites et validées module par module, dans l’ordre suivant :
\begin{enumerate}
    \item Conception du schéma de base de données et des tables nécessaires.
    \item Développement de l’authentification et de la gestion des sessions.
    \item Implémentation du questionnaire et de l’algorithme de calcul de compatibilité de filière (algorithme de scoring).
    \item Création du tableau de bord étudiant et du panel administrateur.
    \item Intégration du module d’intelligence artificielle (chatbot Mistral).
    \item Finalisation des interfaces graphiques et tests fonctionnels.
\end{enumerate}

Le diagramme de Gantt présenté dans la figure~\ref{fig:gantt} décrit le planning du déroulement du projet dans sa globalité, réparti sur les différentes semaines de travail.

\begin{figure}[H]
\centering
\resizebox{\textwidth}{!}{
\begin{ganttchart}[
    x unit=0.15cm,
    y unit title=0.7cm,
    y unit chart=0.8cm,
    vgrid={*4{dotted}, *1{dashed}},
    hgrid,
    time slot format=isodate,
    time slot unit=day,
    title/.append style={fill=brownmain!20},
    title label font=\bfseries\footnotesize,
    bar/.append style={fill=brownlight, draw=brownmain},
    bar label font=\footnotesize,
    milestone/.append style={fill=brownmain, draw=brownmain}
]{2026-02-01}{2026-05-20}
\gantttitlecalendar{year, month} \\
\ganttbar{Analyse \& conception}{2026-02-01}{2026-02-28} \\
\ganttbar{Conception BD \& architecture}{2026-02-01}{2026-02-28} \\
\ganttbar{Développement backend}{2026-03-01}{2026-03-31} \\
\ganttbar{Développement frontend}{2026-03-01}{2026-03-31} \\
\ganttbar{Algorithme d'orientation}{2026-03-01}{2026-03-15} \\
\ganttbar{Intégration IA}{2026-04-01}{2026-04-30} \\
\ganttbar{Tests \& validation}{2026-04-01}{2026-04-30} \\
\ganttbar{Documentation}{2026-05-01}{2026-05-20} \\
\ganttmilestone{Soutenance}{2026-05-20}
\end{ganttchart}
}
\caption{Diagramme de Gantt illustrant le planning du projet}
\label{fig:gantt}
\end{figure}

\subsection{Approche Algorithmique}

L'algorithme de calcul de compatibilité de filière ne se contente pas d'additionner des points. Il repose sur une méthodologie en quatre étapes successives :

\begin{enumerate}

    \item \textbf{Inversion d'échelle :}

    Le questionnaire comporte des affirmations formulées dans les deux sens : certaines sont positives (par exemple : \textit{\og J'aime résoudre des problèmes informatiques\fg{}}), tandis que d'autres sont négatives (par exemple : \textit{\og J'évite de travailler sur des ordinateurs\fg{}}).

    Afin de garantir une interprétation homogène des réponses, les items négativement formulés sont inversés. Cette opération permet de faire en sorte qu'un score élevé traduise systématiquement une forte compatibilité avec la dimension évaluée, indépendamment de la formulation de la question.

    Étant donné que l'échelle de réponse utilisée est comprise entre 1 et 7, l'inversion est réalisée selon la formule :
    \[
        r' = 8 - r
    \]

    où :
    \begin{itemize}
        \item $r$ représente la réponse brute de l'étudiant ;
        \item $r'$ représente la réponse inversée.
    \end{itemize}

    Ainsi :
    \begin{itemize}
        \item une réponse de \textbf{7} devient \textbf{1} ;
        \item une réponse de \textbf{6} devient \textbf{2} ;
        \item une réponse de \textbf{4} reste inchangée.
    \end{itemize}

    Cette transformation garantit que toutes les réponses sont orientées dans le même sens avant l'étape de normalisation statistique.

    \item \textbf{Normalisation par Z-Score :}

    Deux étudiants peuvent avoir des habitudes de notation très différentes : l'un répond systématiquement entre 4 et 5 (étudiant \textit{\og généreux\fg{}}), tandis qu'un autre répond entre 1 et 2 (étudiant \textit{\og sévère\fg{}}). Sans correction, le premier obtiendrait des scores plus élevés, non pas parce qu'il est plus compatible, mais parce qu'il note plus fort. Le Z-Score neutralise ce biais en ramenant les réponses de chaque étudiant à une échelle commune selon la formule :
    \[
        z = \frac{r - \bar{r}}{\sigma}
    \]

    où :
    \begin{itemize}
        \item $r$ : la réponse brute de l'étudiant à une question ;
        \item $\bar{r}$ : la moyenne de toutes ses réponses ;
        \item $\sigma$ : l'écart-type de ses réponses ;
        \item $z$ : le score normalisé final.
    \end{itemize}

    Ainsi, ce qui compte n'est plus la note brute de la réponse, mais \textbf{l'écart par rapport au comportement habituel de l'étudiant}.

    \item \textbf{Agrégation dimensionnelle :}

    Les questions du questionnaire sont regroupées en 8 dimensions comportementales, chacune mesurant une aptitude spécifique liée aux filières de l'établissement. Le tableau~\ref{tab:dimensions} présente la description détaillée de ces 8 dimensions.

\begin{table}[H]
\centering
\begin{tabular}{cp{11cm}}
\toprule
\rowcolor{brownmain!20}
\textbf{Dimension} & \textbf{Aptitude mesurée} \\
\midrule
\textbf{D1} & Capacité à raisonner de manière logique et à résoudre 
des problèmes mathématiques complexes. \\
\textbf{D2} & Affinité avec la programmation, les algorithmes et les 
technologies de l'intelligence artificielle. \\
\textbf{D3} & Intérêt pour la communication numérique, les 
infrastructures réseau et les protocoles de télécommunication. \\
\textbf{D4} & Passion pour les systèmes mécaniques intelligents, 
la robotique et l'automatisation industrielle. \\
\textbf{D5} & Sensibilité aux procédés de fabrication, à la chaîne 
de production et à l'optimisation industrielle. \\
\textbf{D6} & Aptitude à coordonner des équipes, planifier des projets 
et organiser des ressources humaines. \\
\textbf{D7} & Goût pour l'expérimentation scientifique, la modélisation 
théorique et la rédaction de rapports techniques. \\
\textbf{D8} & Conscience des risques systémiques, de la cybersécurité 
et des contraintes environnementales. \\
\bottomrule
\end{tabular}
\caption{Description des 8 dimensions comportementales du questionnaire}
\label{tab:dimensions}
\end{table}
La Table~\ref{tab:dimensions} présente les huit dimensions comportementales évaluées par notre questionnaire. Ce référentiel permet de dresser un profil multidimensionnel de l'étudiant en couvrant l'ensemble des aptitudes requises par les filières de l'ENSA Safi, de la logique pure aux enjeux technologiques et environnementaux.

Pour chaque dimension, on calcule la moyenne des z-scores des questions
qui lui appartiennent. On obtient ainsi un vecteur de 8 scores
dimensionnels $\vec{z} = (z_{D1}, z_{D2}, \ldots, z_{D8})$
qui décrit le profil comportemental complet de l'étudiant,
indépendamment de ses habitudes de notation.

\item \textbf{Transformation Softmax :} Chaque filière est définie par un vecteur de poids sur les 8 dimensions (D1 à D8), traduisant l'importance de chaque aptitude pour cette filière. Ces poids sont répertoriés dans la Matrice de pondération dimensionnelle :

Le tableau~\ref{tab:ponderation} montre la matrice de pondération pour chaque filière.

\begin{table}[ht!]
\centering
\small
\begin{tabularx}{\linewidth}{X c c c c c c c c}
\toprule
\rowcolor{brownmain!20} \textbf{Filière} & \textbf{D1} & \textbf{D2} & \textbf{D3} & \textbf{D4} & \textbf{D5} & \textbf{D6} & \textbf{D7} & \textbf{D8} \\
\midrule
\rowcolor{brownlight} GIIA   & 1.0 & 1.5 & 0.2 & 0.3 & 0.1 & 0.2 & 0.5 & 0.4 \\
GTR    & 0.5 & 0.8 & 1.5 & 0.2 & 0.1 & 0.2 & 0.4 & 0.6 \\
\rowcolor{brownlight} GINDUS & 0.3 & 0.4 & 0.1 & 0.5 & 1.5 & 1.2 & 0.3 & 0.6 \\
GMSI   & 0.8 & 0.5 & 0.4 & 1.5 & 0.3 & 0.2 & 0.4 & 0.5 \\
\rowcolor{brownlight} GATE   & 0.6 & 0.3 & 0.2 & 1.2 & 0.5 & 0.3 & 0.8 & 1.5 \\
GPMA   & 0.4 & 0.2 & 0.1 & 0.3 & 1.2 & 0.4 & 1.5 & 0.6 \\
\bottomrule
\end{tabularx}
\caption{Matrice de pondération dimensionnelle par filière}
\label{tab:ponderation}
\end{table}
La Table~\ref{tab:ponderation} détaille les coefficients de pondération appliqués à chaque dimension pour les six filières de l'ENSA Safi. Ces poids permettent à l'algorithme de calculer un score d'adéquation précis, en privilégiant les aptitudes spécifiques à chaque spécialité.

Pour déterminer la compatibilité de chaque filière, le calcul se déroule en deux étapes successives :
\begin{itemize}
    \item \textbf{Calcul du score brut $S_i$ :} Pour chaque filière $i$, on effectue le produit scalaire entre son vecteur de poids dimensionnels et le profil comportemental de l'étudiant, afin d'obtenir un score reflétant l'adéquation globale entre les aptitudes de l'étudiant et les exigences de la filière :
    \begin{equation}
        S_i = \sum_{k=1}^{8} w_{i,k} \cdot x_k
    \end{equation}
    où :
    \begin{itemize}
        \item $w_{i,k}$ : le poids de la filière $i$ sur la dimension $k$, tel que défini dans la matrice de pondération ci-dessus,
        \item $x_k$ : le score z-normalisé moyen de l'étudiant sur la dimension $k$,
        \item $k \in \{1, \ldots, 8\}$ : l'indice de la dimension comportementale.
    \end{itemize}

    \item \textbf{Normalisation par Softmax en probabilités $P(f_i)$ :} Les scores bruts $S_i$ obtenus ne sont pas directement comparables entre eux, car ils peuvent être négatifs ou de grandeurs très différentes. La fonction Softmax les convertit en probabilités de compatibilité comprises entre 0\% et 100\%, dont la somme vaut toujours 100\%, permettant une lecture claire et équitable du classement des filières :
    \begin{equation}
        P(f_i) = \frac{e^{S_i}}{\sum_{j=1}^{n} e^{S_j}}
    \end{equation}
    où :
    \begin{itemize}
        \item $P(f_i)$ : la probabilité de compatibilité de l'étudiant avec la filière $f_i$,
        \item $S_i$ : le score brut de la filière $i$, calculé par la formule précédente,
        \item $n$ : le nombre total de filières (ici $n = 6$),
        \item $e^{S_j}$ : l'exponentielle du score brut de la filière $j$, assurant que toutes les valeurs sont positives avant normalisation.
    \end{itemize}
\end{itemize}

\end{enumerate}
\subsection{Architecture logicielle (MVC)}

Pour garantir une maintenance aisée et une séparation claire des responsabilités, l'application suit l'architecture logicielle \textbf{Model-View-Controller (MVC)}. La Table~\ref{tab:mvc} présente la répartition des différents fichiers et technologies du projet selon ces trois couches.

\begin{table}[H]
\centering
\renewcommand{\arraystretch}{1.4}
\begin{tabularx}{\linewidth}{L{2.5cm} L{4cm} X}
\toprule
\rowcolor{brownmain!20} \textbf{Couche} & \textbf{Technologie} & \textbf{Rôle} \\
\midrule
\rowcolor{brownlight} \textbf{Modèle (M)} & SQLite3 & Stockage persistant des données (étudiants, résultats, questions) \\
\textbf{Vue (V)} & HTML5, CSS3, JS & Affichage des pages et interfaces utilisateur \\
\rowcolor{brownlight} \textbf{Contrôleur (C)} & Python / Flask & Traitement des requêtes, logique métier, routage \\
\textbf{Service IA} & API Mistral & Assistant conversationnel en langage naturel \\
\bottomrule
\end{tabularx}
\caption{Couches de l'architecture MVC du projet}
\label{tab:mvc}
\end{table}

Cette organisation permet de découpler la logique métier (Flask) de l'interface utilisateur (Jinja2/CSS) et de la gestion des données (SQLite). Une telle structure facilite l'évolution du système, comme l'ajout futur de nouvelles fonctionnalités ou le remplacement de la base de données sans impacter l'affichage.
\subsection{Organisation du code}

Le projet est structuré selon les répertoires suivants :

\begin{lstlisting}[style=filetree, caption={Arborescence du projet}, captionpos=b]
mini-projet/
|-- backend/
|   |-- app.py               # Entree Flask
|   |-- models/
|   |   |-- database.py      # SQLite
|   |-- routes/
|   |   |-- main_routes.py   # Pages HTML
|   |   |-- api_routes.py    # API REST
|   |   |-- admin_routes.py  # Administration
|   |-- services/
|   |   |-- ai_service.py    # Mistral AI
|   |  
|   |-- utils/
|       |-- helpers.py       # Decorateurs Jinja2
|       |-- config.py        # Configue globale
|-- database/
|   |-- schema.sql           # Tables SQL
|   |-- data.sql             # Donnees initiales
|   |-- ensa_orientation.db  # Fichier SQLite
|-- frontend/
    |-- html/                # Templates Jinja2
    |-- css/                 # Styles CSS
    |-- js/                  # Scripts JS
\end{lstlisting}

\subsection{Diagramme de flux général}
La figure~\ref{fig:flux_global} illustre le chemin emprunté par une requête depuis le navigateur de l’étudiant jusqu’à la réponse du serveur.

% --- Définition des couleurs pour la figure ---
\definecolor{ensablue}{HTML}{6B3A2A}   % brun foncé  (remplace bleu marine)
\definecolor{ensalight}{HTML}{FDF6F0}  % crème chaud (remplace gris clair)
\definecolor{accentgray}{HTML}{A0785A} % brun moyen  (remplace gris neutre)

\begin{figure}[ht!]
\centering
\resizebox{\textwidth}{!}{%
\begin{tikzpicture}[
  node distance=1cm,
  box/.style={draw=ensablue, fill=ensalight, rounded corners=4pt, 
              minimum width=2.8cm, minimum height=0.9cm, font=\small\bfseries, align=center},
  db/.style={draw=ensablue, fill=white, rounded corners=2pt, 
             minimum width=2.8cm, minimum height=0.9cm, font=\small, align=center},
  arr/.style={-Stealth, thick, color=ensablue},
  note/.style={font=\scriptsize\color{accentgray}, align=center}
]
  % Placement des nœuds avec espacement réduit
  \node[box] (browser) {Navigateur\\étudiant};
  \node[box, right=1.4cm of browser] (flask) {Serveur Flask\\(app.py)};
  \node[box, right=1.4cm of flask]   (routes) {Blueprints\\(routes)};
  \node[db,  below=1.2cm of routes]  (db)     {SQLite\\(database.py)};
  \node[db,  right=1.4cm of routes]  (ai)     {Mistral AI\\(api externe)};

  % Tracé des flux
  \draw[arr] (browser) -- node[above,note]{HTTP} (flask);
  \draw[arr] (flask)   -- node[above,note]{routage} (routes);
  \draw[arr] (routes)  -- node[right,note]{CRUD} (db);
  \draw[arr] (routes)  -- node[above,note]{/api/chat} (ai);
  \draw[arr] (db)      -- (routes);
  
  % Flux de retour
  \draw[arr, dashed] (ai) to[bend left=20] node[below,note]{JSON} (routes);
  \draw[arr, dashed] (routes) to[bend right=35] node[below,note]{HTML/JSON} (browser);

\end{tikzpicture}%
}
\caption{Flux général des données dans ENSA Guide}
\label{fig:flux_global}
\end{figure}
La Figure 2 illustre le cycle de vie d'une requête au sein de la plateforme. On y observe le rôle central du serveur Flask qui orchestre les échanges : il réceptionne les requêtes HTTP du navigateur, interroge la base de données SQLite via les Blueprints pour la persistance, et sollicite l'API externe Mistral AI pour enrichir l'expérience utilisateur par une assistance intelligente en temps réel.

% ============================================================
\clearpage
\section{Outils et technologies}
% ============================================================

\subsection{Backend : Python et Flask}
\begin{minipage}[c]{0.75\linewidth}
Flask \cite{flask} est utilisé comme micro-framework web. Son mécanisme de \textit{Blueprints} permet
de découper logiquement le code en modules indépendants (\texttt{main\_routes}, \texttt{api\_routes},
\texttt{admin\_routes}). La bibliothèque \textbf{NumPy} \cite{numpy} est employée pour les calculs vectoriels
du test d'orientation (normalisation par z-scores, softmax). La sécurité des mots de passe est
assurée par \texttt{werkzeug.security} \cite{werkzeug} avec hachage \textit{scrypt}. Le langage Python \cite{python} apporte une flexibilité remarquable.
\end{minipage}%
\hfill
\begin{minipage}[c]{0.15\linewidth}
    \centering
    \includegraphics[width=0.8\linewidth]{images/python.png}\\[0.2cm]
    \includegraphics[width=0.8\linewidth]{images/logo_flask_python.png}
\end{minipage}

\vspace{0.4cm}
\subsection{Base de données : SQLite3}
\begin{minipage}[c]{0.75\linewidth}
SQLite3 \cite{sqlite} est retenu pour sa légèreté et son intégration native en Python, sans serveur externe
à configurer. Le fichier \texttt{ensa\_orientation.db} est généré automatiquement au démarrage
si absent, à partir du schéma (\texttt{schema.sql}) et des données initiales (\texttt{data.sql}).
\end{minipage}%
\hfill
\begin{minipage}[c]{0.15\linewidth}
    \centering
    \includegraphics[width=\linewidth]{images/logo_sqlite.png}
\end{minipage}

\vspace{0.4cm}
\subsection{Frontend : HTML5, CSS3 et JavaScript}
\begin{minipage}[c]{0.75\linewidth}
L'interface est développée en HTML/CSS/JavaScript \cite{mdn} sans framework tiers. Le moteur de templates
Jinja2 \cite{jinja} permet le rendu dynamique des pages côté serveur. L'API \texttt{Fetch} JavaScript est
utilisée pour les interactions asynchrones (soumission du test, chatbot).
\end{minipage}%
\hfill
\begin{minipage}[c]{0.15\linewidth}
    \centering
    \includegraphics[width=0.9\linewidth]{images/logo_html_css_js.png}
\end{minipage}

\vspace{0.4cm}
\subsection{Intelligence artificielle : Mistral AI}
\begin{minipage}[c]{0.75\linewidth}
Le chatbot est connecté à l'API Mistral AI \cite{mistral} via le service dédié \texttt{ai\_service.py},
en utilisant le modèle \textbf{\texttt{mistral-tiny}}. Chaque requête est précédée d'un
\textit{system prompt} définissant le rôle de l'assistant : orientation exclusive vers les
six filières de l'ENSA Safi, réponses concises (2 à 3 phrases), sans formatage Markdown,
avec empathie. En cas d'indisponibilité de l'API (timeout, erreur HTTP), un message de
repli explicite est retourné à l'utilisateur, garantissant la continuité de l'expérience.
\end{minipage}%
\hfill
\begin{minipage}[c]{0.15\linewidth}
    \centering
    \includegraphics[width=\linewidth]{images/logo_mistral.png}
\end{minipage}
% ============================================================
\clearpage
\section{Implémentation}
% ============================================================


\subsection{Schéma de la base de données}

La base de données comporte quatre tables principales, définies dans \texttt{schema.sql} :

% UPDATED: caption clarified — inline comments are annotations for readability, not present in schema.sql
\begin{lstlisting}[language=SQL]
CREATE TABLE specialties (
    id              INTEGER PRIMARY KEY AUTOINCREMENT,
    code            TEXT    NOT NULL UNIQUE,
    name            TEXT    NOT NULL,
    description     TEXT    NOT NULL,
    career_paths    TEXT,
    subjects        TEXT,
    additional_info TEXT    -- JSON : programme par annee et semestre
);

CREATE TABLE questions (
    id         INTEGER PRIMARY KEY AUTOINCREMENT,
    texte      TEXT    NOT NULL,
    texte_en   TEXT,             -- traduction anglaise
    texte_ar   TEXT,             -- traduction arabe
    dimensions TEXT    NOT NULL  -- ex. 'D1,D3'
);

CREATE TABLE users (
    id               INTEGER PRIMARY KEY AUTOINCREMENT,
    email            TEXT    UNIQUE NOT NULL,
    access_code      TEXT    NOT NULL,   -- hache via werkzeug scrypt
    nom_complet      TEXT,
    role             TEXT,               -- 'student' ou 'admin'
    is_active        INTEGER DEFAULT 1,
    last_result_json TEXT,
    profile_type     TEXT,
    date_test        DATETIME
);
\end{lstlisting}

\begin{lstlisting}[language=SQL, caption={Schéma SQL complet (schema.sql)}, captionpos=b, firstnumber=30]
CREATE TABLE results (
    id           INTEGER PRIMARY KEY AUTOINCREMENT,
    user_id      INTEGER,
    date_test    DATETIME,
    result_json  TEXT,
    profile_type TEXT,
    duration     INTEGER,          -- duree du test en secondes
    FOREIGN KEY(user_id) REFERENCES users(id)
);
\end{lstlisting}
Le champ \texttt{additional\_info} de \texttt{specialties} contient un JSON du programme sur trois ans.
\subsection{Initialisation de la base de données}
% UPDATED: replaced placeholder ellipsis with actual path resolution pattern
\begin{lstlisting}[language=Python, caption={Initialisation de la base de données (database.py)},captionpos=b]
import sqlite3, os
ROOT = os.path.dirname(os.path.dirname(os.path.dirname(__file__)))
DB_PATH = os.path.join(ROOT, 'database', 'ensa_orientation.db')
def get_db_connection():
    conn = sqlite3.connect(DB_PATH)
    conn.row_factory = sqlite3.Row
    return conn
def init_db():
    if not os.path.exists(DB_PATH):
        conn = get_db_connection()
        try:
            schema_path = os.path.join(ROOT, 'database', 'schema.sql')
            with open(schema_path, 'r', encoding='utf-8') as f:
                conn.executescript(f.read())
            data_path = os.path.join(ROOT, 'database', 'data.sql')
            if os.path.exists(data_path):
                with open(data_path, 'r', encoding='utf-8') as f:
                    conn.executescript(f.read())
        except Exception as e:
            print(f"Schema error: {e}")
        finally:
            conn.close()
\end{lstlisting}
L'initialisation est déclenchée au démarrage du serveur via \texttt{init\_db}.
Si le fichier SQLite est absent, le schéma et les données sont chargés automatiquement.
\noindent\textbf{Remarque :} ce mécanisme crée un problème connu~: si la base de données
existe déjà, les modifications apportées à \texttt{data.sql} ne sont \emph{pas} appliquées
automatiquement. L'utilisateur doit supprimer manuellement le fichier \texttt{.db} pour
forcer la réinitialisation.

\subsection{Initialisation de l'application Flask}

Le fichier \texttt{app.py} adopte le patron \textit{Application Factory} via \texttt{create\_app},
isolant la configuration Flask pour faciliter les tests et le déploiement multi-environnement.
Il définit les chemins absolus des templates Jinja2 et fichiers statiques, enregistre les trois
Blueprints, déclare un filtre \texttt{from\_json} pour le décodage JSON dans les templates,
expose une route \texttt{/assets/<filename>} pour les ressources graphiques, et injecte
automatiquement \texttt{user}, \texttt{logged\_in} et \texttt{role} dans tous les templates
via \texttt{inject\_user}.
\begin{lstlisting}[language=Python, caption={Fabrique d'application et enregistrement des Blueprints (app.py)},captionpos=b]
def create_app():
    base_dir     = os.path.dirname(os.path.dirname(os.path.abspath(__file__)))
    template_dir = os.path.join(base_dir, 'frontend', 'html')
    static_dir   = os.path.join(base_dir, 'frontend')

    app = Flask(__name__,
                template_folder=template_dir,
                static_folder=static_dir,
                static_url_path='/static')

    app.secret_key = 'ensa_safi_secret_key'
    app.config['PERMANENT_SESSION_LIFETIME'] = timedelta(days=7)
    app.config['SESSION_COOKIE_HTTPONLY']    = True
    app.config['SESSION_COOKIE_SAMESITE']    = 'Lax'

    init_db()
    app.template_filter('from_json')(from_json_filter)
    app.register_blueprint(main_bp)
    app.register_blueprint(admin_bp)
    app.register_blueprint(api_bp)
    return app
\end{lstlisting}



\subsection{Authentification et gestion des sessions}

La route \texttt{/login} (méthode POST) implémente un flux d'authentification
séquentiel. Avant toute interrogation de la base, l'adresse e-mail est validée
pour s'assurer qu'elle appartient au domaine académique (\texttt{@uca.ac.ma} ou
\texttt{@ensas.uca.ma}). La vérification du code d'accès est ensuite réalisée
par \texttt{check\_password\_hash} de \texttt{werkzeug.security}, qui supporte
nativement l'algorithme \textit{scrypt}. Un compte désactivé (\texttt{is\_active = 0})
est rejeté avec un message explicite, indépendamment de la validité du code.
\begin{lstlisting}[language=Python, caption={Flux d'authentification simplifié (main\_routes.py)}, captionpos=b]
if not email.endswith('@uca.ac.ma') and \\
   not email.endswith('ensas.uca.ma'):
    return render_template('login.html',
                           error="Email academique requis")

user = conn.execute(
    'SELECT * FROM users WHERE email = ?', (email,)
).fetchone()

if user and check_password_hash(user['access_code'], code):
    if user['is_active'] == 0:
        flash('Compte desactive par l\'administrateur.', 'error')
        return render_template('login.html')
    session.clear()
    session['user_id']   = user['id']
    session['user_name'] = user['nom_complet'] or email.split('@')[0]
    session['role']      = user['role'] or 'student'
    if session['role'] == 'admin':
        return redirect(url_for('admin_routes.admin_analytics'))
\end{lstlisting}

Un mécanisme de \textit{fusion de session} complète ce flux~: si l'utilisateur a
effectué le test d'orientation \emph{avant} de se connecter, ses résultats temporaires
--- stockés dans \texttt{session['temp\_results']} --- sont automatiquement persistés
en base de données dès l'authentification réussie. Ce comportement garantit qu'aucune
donnée n'est perdue lors d'une connexion tardive.

\subsection{Rendu dynamique des fiches de spécialité}

La route \texttt{/filiere/<code\_filiere>} récupère la spécialité par son code dans
la table \texttt{specialties} et effectue un mappage explicite vers les noms de champs
attendus par le template Jinja2 :
\newpage
\begin{lstlisting}[language=Python, caption={Route de rendu des fiches de spécialité (main\_routes.py)},captionpos=b]
@main_bp.route('/filiere/<code_filiere>')
def filiere_detail(code_filiere):
    conn = get_db_connection()
    filiere = conn.execute(
        'SELECT * FROM specialties WHERE code = ?',
        (code_filiere,)
    ).fetchone()
    conn.close()
    if filiere:
        mapped_filiere = {
            'code':           filiere['code'],
            'nom':            filiere['name'],
            'description':    filiere['description'],
            'debouches':      filiere['career_paths'],
            'mots_cles':      filiere['subjects'],
            'formation_json': filiere['additional_info']
        }
        return render_template('filiere_detail.html', filiere=mapped_filiere)
    return "Filiere non trouvee", 404
\end{lstlisting}

Ce mappage constitue une couche d'adaptation entre le schéma SQL anglophone adopté lors
de la migration et les templates Jinja2 initialement conçus avec des noms de champs
francophones. Il isole les templates de toute modification ultérieure du schéma,
conformément au principe de découplage entre la présentation et la persistance des données.

\subsection{Calcul du score d'orientation}

Le test soumet les réponses en format JSON à la route \texttt{/api/submit-test}.
L'algorithme de scoring se déroule en quatre étapes :

\begin{enumerate}
    \item \textbf{Z-normalisation des réponses} : chaque réponse brute est normalisée
          par rapport à la moyenne et à l'écart-type des réponses de l'utilisateur,
          afin de neutraliser les biais de tendance centrale.
    \item \textbf{Calcul des scores dimensionnels} : les questions sont associées à
          des dimensions (D1 à D8, représentant des aptitudes telles que la logique,
          l'informatique, les réseaux, etc.). Le score de chaque dimension est la
          moyenne des z-scores des questions qui lui correspondent.
    \item \textbf{Pondération par filière} : chaque filière est définie par un vecteur
          de poids sur les dimensions, configuré dans \texttt{config.py}. Le score brut
          d'une filière est le produit scalaire de ce vecteur et des scores dimensionnels.
    \item \textbf{Softmax} : les scores bruts sont convertis en probabilités de
          compatibilité (0--100\%) via la fonction softmax, ce qui garantit que leur
          somme est toujours égale à 100\%.
\end{enumerate}

La formule appliquée est :
\begin{equation}
    P(f) = \frac{e^{S(f)}}{\displaystyle\sum_{k} e^{S(k)}} \times 100
    \quad \text{avec} \quad
    S(f) = \sum_{d \in D} w_{f,d} \cdot \bar{z}_d
\end{equation}

où :
\begin{itemize}
    \item $\boldsymbol{w_{f,d}}$ : le poids de la filière $f$ sur la dimension $d$,
    \item $\boldsymbol{\bar{z}_d}$ : le score z-normalisé moyen de la dimension $d$.
\end{itemize}
Ces poids, définis statiquement dans le fichier \texttt{config.py}, encodent l'affinité dimensionnelle de chaque filière. Ils correspondent aux coefficients détaillés précédemment dans la \textbf{Table~\ref{tab:ponderation}} (page~\pageref{tab:ponderation}), qui servent de base au calcul du score d'adéquation.
\begin{lstlisting}[language=Python, caption={Calcul softmax des scores de compatibilité (api\_routes.py)},captionpos=b]
# Ponderation par filiere
filiere_scores = {}
for f_code, weights in WEIGHTS.items():
    score = sum(weights.get(d, 0) * dim_scores[d]
                for d in DIMENSIONS)
    filiere_scores[f_code] = score

# Softmax -> probabilites de compatibilite
exp_scores = {f: np.exp(s) for f, s in filiere_scores.items()}
sum_exp = sum(exp_scores.values())
probabilities = {
    f: (exp_scores[f] / sum_exp) * 100
    for f in exp_scores
}
\end{lstlisting}

Une fois les probabilités calculées, un type de profil est déterminé selon trois
cas mutuellement exclusifs et persisté dans la table \texttt{results}~:
\begin{itemize}
    \item \textbf{Spécialiste} \texttt{<code>} : si la filière en tête dépasse 40\,\%
          de compatibilité.
    \item \textbf{Profil Hybride} \texttt{<code1>/<code2>} : si l'écart entre les deux
          premières filières est inférieur à 5 points de pourcentage.
    \item \textbf{Profil Équilibré} : dans tous les autres cas.
\end{itemize}

\subsection{Tableau de bord étudiant}

Le tableau de bord constitue l'espace personnel de l'étudiant authentifié au sein de la plateforme. Accessible via la route \texttt{/dashboard}, protégée par une vérification stricte de session côté serveur, cette interface centralise l'ensemble des données propres à l'utilisateur et lui en restitue une lecture structurée et exploitable.

Lorsqu'un étudiant a effectué plusieurs passages du test d'orientation, la route agrège l'intégralité de l'historique stocké dans la table \texttt{results} et calcule dynamiquement un \textbf{profil global consolidé}. Ce profil est obtenu en moyennant, filière par filière, les scores de compatibilité issus de l'ensemble des tests réalisés, puis en appliquant les mêmes règles de classification que pour un test individuel --- à savoir la détermination d'un profil Spécialiste, Hybride ou Équilibré selon les seuils définis. 

Afin de mieux visualiser ce processus de traitement des données historiques, le schéma suivant détaille le mécanisme d'agrégation mis en œuvre:
\begin{figure}[htbp]
\centering
\begin{tikzpicture}[
  every node/.style={font=\small},
  bloc/.style={draw=brownmain, fill=brownlight, rounded corners=5pt,
               minimum width=3.2cm, minimum height=1.1cm, align=center, line width=0.8pt},
  procstep/.style={draw=brownmain, fill=white, rounded corners=5pt,
               minimum width=5.5cm, minimum height=1.1cm, align=center, line width=0.8pt},
  arr/.style={-Stealth, thick, color=brownmain},
  lbl/.style={font=\scriptsize\color{brownmid}, align=center}
]

%% --- Ligne 1 : historique tests (horizontalement centrée) ---
\node[bloc] (t1) at (-4, 0) {\textbf{Test 1} \\ \scriptsize résultats JSON};
\node[bloc] (t2) at ( 0, 0) {\textbf{Test 2} \\ \scriptsize résultats JSON};
\node[bloc] (tn) at ( 4, 0) {\textbf{Test N} \\ \scriptsize résultats JSON};

%% --- Encadrement manuel des Tests ---
\draw[draw=brownmid, dashed, rounded corners=6pt, line width=0.7pt]
  (-6, 0.9) rectangle (6, -0.9);
\node[lbl] at (0, 1.1) {Historique étudiant — table \texttt{results}};

%% --- Ligne 2 : agrégation ---
\node[procstep] (agg) at (0, -2.5)
  {\textbf{Agrégation} \\ \scriptsize filière par filière (moyenne des scores)};

%% --- Ligne 3 : profil global ---
\node[procstep] (profil) at (0, -4.7)
  {\textbf{Profil global} \\ \scriptsize consolidé};

%% --- Ligne 4 : classification ---
\node[bloc, minimum width=6cm] (classif) at (0, -6.9)
  {\textbf{Classification} \\ \scriptsize Spécialiste / Hybride / Équilibré};

%% --- Flèches vers agrégation ---
\draw[arr] (t1.south) -- ++(0,-0.6) -| (agg.north);
\draw[arr] (t2.south) -- (agg.north);
\draw[arr] (tn.south) -- ++(0,-0.6) -| (agg.north);

%% --- Flèches descendantes principales ---
\draw[arr] (agg)    -- node[right, lbl]{calcul du score moyen} (profil);
\draw[arr] (profil) -- node[right, lbl]{seuils : 40\% et 5~pts} (classif);

\end{tikzpicture}
\caption{Flux d'agrégation des résultats dans le tableau de bord étudiant}
\label{fig:dashboard_flow}
\vspace{0.3cm}
\end{figure}

Cette analyse permet d'obtenir des résultats plus fiables. En réduisant l'effet des variations temporaires (comme l'humeur ou les conditions du test), elle met en évidence les véritables préférences de l'étudiant. Ainsi, l'étudiant peut identifier une orientation plus stable et suivre l'évolution de son profil dans le temps.
\newpage

\subsection{Module d'administration}

Le module d'administration offre à l'administrateur de la plateforme un environnement
de supervision complet, articulé autour de deux routes distinctes, toutes deux protégées
par le décorateur \texttt{@admin\_required}, qui contrôle l'accès au niveau de chaque
requête HTTP et redirige systématiquement tout utilisateur non habilité.

\textbf{Page analytique (\texttt{/admin/analytics})}~: cette vue interroge l'ensemble
de la table \texttt{results} afin de produire trois indicateurs statistiques injectés
dans le template via le dictionnaire \texttt{chart\_data}~:
\begin{itemize}
    \item Un \textbf{histogramme} représentant la répartition des recommandations
          par filière, permettant d'identifier les orientations les plus fréquemment
          attribuées au sein de la promotion.
    \item Un \textbf{diagramme en anneau} illustrant la distribution des profils
          types (Spécialiste, Hybride, Équilibré) et reflétant la diversité des
          configurations cognitives et motivationnelles de la population étudiante.
    \item Une \textbf{courbe de tendance mensuelle} visualisant l'évolution du volume
          de tests réalisés dans le temps, utile pour évaluer l'adoption de la plateforme.
\end{itemize}

La figure~\ref{fig:admin_architecture} résume l'architecture globale de ce module d'administration. Afin d'assurer une gestion rigoureuse de la plateforme, ce module centralise le pilotage des données et la supervision des utilisateurs via une interface sécurisée, reposant sur des routes protégées et une API REST dédiée.

\begin{figure}[htbp]
\centering
\begin{tikzpicture}[
  every node/.style={font=\small},
  garde/.style={draw=brownmain, fill=brownmain!20, rounded corners=5pt,
                minimum width=6cm, minimum height=1cm, align=center, line width=0.8pt},
  route/.style={draw=brownmain, fill=brownlight, rounded corners=5pt,
                minimum width=5cm, minimum height=1cm, align=center, line width=0.8pt},
  kpi/.style={draw=brownmid, fill=white, rounded corners=4pt,
              minimum width=4.5cm, minimum height=0.9cm, align=center, line width=0.7pt},
  api/.style={draw=brownmid, fill=brownlight!60, rounded corners=4pt,
              minimum width=4.5cm, minimum height=0.9cm, align=center, line width=0.7pt},
  arr/.style={-Stealth, thick, color=brownmain},
]

\node[garde] (guard) at (5, 0)
  {\texttt{@admin\_required} \\ \scriptsize contrôle d'accès HTTP};

% Routes
\node[route] (analytics) at (1.5, -2) {\texttt{/admin/analytics} \\ \scriptsize Page analytique};
\node[route] (users)     at (8.5, -2) {\texttt{/admin/users} \\ \scriptsize Gestion des comptes};

\draw[arr] (guard.south) -- ++(0,-0.5) -| (analytics.north);
\draw[arr] (guard.south) -- ++(0,-0.5) -| (users.north);

% Analytics children
\node[kpi] (bar)   at (4, -3.8) {Histogramme \\ \scriptsize répartition par filière};
\node[kpi] (donut) at (4, -5.2) {Diagramme anneau \\ \scriptsize distribution profils};
\node[kpi] (trend) at (4, -6.6) {Courbe de tendance \\ \scriptsize volume mensuel};
\node[api] (csv)   at (4, -8.0) {\texttt{GET} export\_csv \\ \scriptsize données brutes CSV};

\draw[thick, color=brownmain] (analytics.south) -- (1.5, -8.0);
\draw[arr] (1.5, -3.8) -- (bar.west);
\draw[arr] (1.5, -5.2) -- (donut.west);
\draw[arr] (1.5, -6.6) -- (trend.west);
\draw[arr] (1.5, -8.0) -- (csv.west);

% Users children
\node[api] (add)    at (11, -3.8) {\texttt{POST} add\_student \\ \scriptsize Créer un compte};
\node[api] (toggle) at (11, -5.2) {\texttt{POST} toggle\_user \\ \scriptsize Activer / désactiver};
\node[api] (del)    at (11, -6.6) {\texttt{DELETE} delete\_user \\ \scriptsize Supprimer compte};

\draw[thick, color=brownmain] (users.south) -- (8.5, -6.6);
\draw[arr] (8.5, -3.8) -- (add.west);
\draw[arr] (8.5, -5.2) -- (toggle.west);
\draw[arr] (8.5, -6.6) -- (del.west);

\end{tikzpicture}
\caption{Architecture du module d'administration : routes protégées, indicateurs et API REST}
\label{fig:admin_architecture}
\vspace{0.3cm}
\end{figure}

Comme l'illustre la Figure~\ref{fig:admin_architecture}, le module est scindé en deux piliers fonctionnels : le volet analytique (KPIs) et la gestion des utilisateurs (\texttt{/admin/users}). Cette interface permet de suivre dynamiquement le nombre de tests effectués et la filière conseillée pour chaque profil.

L'administrateur dispose d'un contrôle total via une API REST dédiée :

{\small
\begin{itemize}[leftmargin=1cm, labelsep=0.2cm, nosep]
    \item \textbf{Création} d'un nouveau compte étudiant (\texttt{POST /api/add\_student}).
    \item \textbf{Activation/Désactivation} d'un compte (\texttt{POST /api/toggle\_user/<id>}).
    \item \textbf{Réinitialisation} du code d'accès (\texttt{POST /api/update\_password}).
    \item \textbf{Suppression} entraînant la purge de l'historique (\texttt{DELETE /api/delete\_user/<id>}).
\end{itemize}
}

\vspace{0.3cm}
\noindent
\begin{infobox}
\textbf{Sécurité :} Une garde logicielle explicite bloque toute tentative de suppression d'un compte disposant du rôle \texttt{admin}, préservant ainsi l'intégrité de l'accès à la plateforme.
\end{infobox}
\subsection{Service d'intelligence artificielle}

Le service d'assistance conversationnelle constitue l'un des éléments différenciateurs
de la plateforme. Implémenté au sein du module \texttt{ai\_service.py}, il expose un
assistant virtuel spécialisé, directement accessible depuis un widget flottant présent
sur l'ensemble des pages de l'application. Son objectif est de permettre à l'étudiant
d'obtenir des réponses ciblées à ses interrogations d'orientation, en langage naturel,
sans nécessiter de navigation supplémentaire.

Le service repose sur l'API Mistral AI, interrogée via des requêtes HTTP POST
adressées à l'endpoint \texttt{https://api.mistral.ai/v1/chat/completions}.
Le modèle \texttt{mistral-tiny} est sélectionné pour son équilibre entre rapidité
de traitement et qualité de génération. 

Chaque requête envoyée à l’assistant est précédée d’un message spécial (\textit{system prompt}) qui définit son rôle. L’assistant doit uniquement répondre sur les six filières de l’ENSA Safi, en 2 à 3 phrases, sans utiliser de formatage.

Ensuite, l’application nettoie automatiquement la réponse pour supprimer les éléments inutiles (comme \texttt{**} ou \texttt{*}), afin d’afficher un texte propre dans l’interface.

Le système gère également les erreurs afin d’assurer un bon fonctionnement :
\begin{itemize}
    \item \textbf{Délai dépassé} : l’utilisateur est invité à réessayer plus tard.
    \item \textbf{Erreur de l’API} : un message indique que le service est temporairement indisponible.
    \item \textbf{Erreur imprévue} : un message simple est affiché sans montrer de détails techniques.
\end{itemize}
La figure~\ref{fig:ai_pipeline} illustre le pipeline complet de traitement de ces requêtes conversationnelles. Ce schéma détaille le flux de données, depuis l'interaction initiale de l'utilisateur jusqu'à la restitution d'une réponse enrichie par l'intelligence artificielle, tout en assurant une continuité de service en cas d'incident technique.

\begin{figure}[htbp]
\centering
\begin{tikzpicture}[
  every node/.style={font=\normalsize},
  stage/.style={draw=brownmain, fill=brownlight, rounded corners=5pt,
                minimum width=5cm, minimum height=1.1cm, align=center, line width=0.8pt},
  ext/.style={draw=brownmid, fill=white, rounded corners=5pt,
              minimum width=5cm, minimum height=1.1cm, align=center, line width=0.9pt},
  arr/.style={-Stealth, thick, color=brownmain},
  note/.style={font=\small\color{brownmid}, align=center},
  err/.style={draw=codered!70, fill=codered!8, rounded corners=5pt,
              minimum width=5.5cm, minimum height=1.2cm, align=center, line width=0.7pt}
]

\node[stage] (widget)   at (0, 0)     {\textbf{Widget flottant} \\ \small message étudiant};
\node[stage] (route)    at (0, -2.0)  {\textbf{Route Flask} \\ \small \texttt{/api/chat}};
\node[ext]   (mistral)  at (0, -4.0)  {\textbf{API Mistral AI} \\ \small \texttt{mistral-tiny}};
\node[stage] (postproc) at (0, -6.0)  {\textbf{Post-traitement} \\ \small nettoyage Markdown};
\node[stage] (html)     at (0, -8.0)  {\textbf{Réponse HTML} \\ \small interface chat};

\draw[arr] (widget) -- node[right, note]{JSON} (route);
\draw[arr] (route) -- node[right, note]{POST HTTP} (mistral);
\draw[arr] (mistral) -- node[right, note]{JSON} (postproc);
\draw[arr] (postproc) -- node[right, note]{texte} (html);

\node[err, anchor=west] (errs) at (3.5, -6.5)
  {\textbf{Gestion des erreurs} \\
   \small Timeout \;·\; HTTPError \\ \small Exception imprévue \\
   $\Rightarrow$ \textbf{repli en français}};
   
\draw[-Stealth, thick, color=codered!70] (mistral.south east) |- (errs.west);
\draw[-Stealth, thick, color=codered!70] (errs.south) |- (html.east);

\end{tikzpicture}
\caption{Pipeline de traitement des requêtes du service d'intelligence artificielle}
\label{fig:ai_pipeline}
\vspace{0.3cm}
\end{figure}
\FloatBarrier

Comme l'illustre ce schéma, le serveur Flask joue le rôle de chef d'orchestre. Il reçoit les messages de l'étudiant, les prépare pour l'API Mistral AI, puis nettoie la réponse reçue pour qu'elle s'affiche correctement dans le chat. De plus, un système de sécurité est en place : si l'API met trop de temps à répondre ou rencontre une erreur, le serveur envoie automatiquement un message de secours. Cela garantit que l'étudiant reçoive toujours une réponse, même simplifiée, au lieu d'un message d'erreur.
% ============================================================
\clearpage
\section{Interface}
% ============================================================

% ----------------------------------------------------------
\subsection{Accueil}
% ----------------------------------------------------------

\textbf{Route :} \texttt{/}

\textbf{Overview :}
Page d'entrée de la plateforme, présentant les statistiques globales et offrant un accès
direct aux six filières de l'ENSA Safi via une grille de cartes interactives.

\textbf{Key Elements :}
\begin{itemize}
    \item Section \textit{hero} avec indicateurs statistiques (tests, filières).
    \item Grille de six cartes cliquables, une par filière (GIIA, GTR, GINDUS, GMSI, GPMA, GATE).
    \item Barre de navigation avec liens vers le test d'orientation, la connexion et le tableau de bord.
    \item Appel à l'action (CTA) invitant l'utilisateur à démarrer le test d'orientation.
\end{itemize}
\textbf{Interface Description :}
La page est structurée en deux blocs verticaux principaux. Le \textit{hero} occupe la partie supérieure
avec un fond dégradé, un titre prominent et trois compteurs dynamiques injectés par Jinja2. En dessous,
une grille CSS à trois colonnes affiche les cartes de filière, chacune composée d'une icône, d'un code
court, d'un intitulé et d'un lien vers la fiche détaillée. La navbar est fixe et adaptative, intégrant
Le statut de connexion de l'utilisateur y est également affiché (bouton \textit{Connexion} ou \textit{Mon espace}). La figure~\ref{fig:accueil_hero} présente la section hero de la page d'accueil, tandis que la figure~\ref{fig:accueil_grille} expose la grille de présentation des filières.

\begin{figure}[H]
\centering
\includegraphics[width=1\textwidth]{images/Image1.png}
\caption{Page d'accueil — section hero}
\label{fig:accueil_hero}
\end{figure}

\begin{figure}[H]
\centering
\includegraphics[width=1\textwidth]{images/Image2.png}
\caption{Page d'accueil — grille des six filières}
\label{fig:accueil_grille}
\end{figure}

% ----------------------------------------------------------
\subsection{Fiche filière}
% ----------------------------------------------------------

\textbf{Route :} \texttt{/filiere/<code>}

\textbf{Overview :}
Page de détail d'une filière, présentant l'intégralité des informations académiques et
professionnelles permettant à l'étudiant de prendre une décision éclairée.

\textbf{Key Elements :}
\begin{itemize}
    \item Description générale et identité de la filière.
    \item Programme structuré par année et semestre.
    \item Débouchés professionnels sous forme de liste.
    \item Mots-clés thématiques (domaines de compétence associés).
\end{itemize}

\textbf{Interface Description :}
La mise en page adopte un layout deux colonnes : une barre latérale (sidebar) à gauche regroupant
les métadonnées de la filière (code, mots-clés, débouchés) et une zone principale à droite
affichant la description et le programme détaillé. Le programme est rendu dynamiquement
à partir du JSON \texttt{additional\_info} via le filtre Jinja2 \texttt{from\_json},
organisé en blocs d'accordéon ou de tableaux par semestre. Un bouton d'appel à l'action
redirige vers le test d'orientation. La figure~\ref{fig:fiche_desc} montre la présentation détaillée du programme et des débouchés de la filière GTR, et la figure~\ref{fig:fiche_comp} affiche les compétences et technologies associées.

\begin{figure}[H]
\centering
\includegraphics[width=1\textwidth]{images/Image3.png}
\caption{Fiche filière — description, programme par semestre et débouchés professionnels du filière GTR}
\label{fig:fiche_desc}
\end{figure}

\begin{figure}[H]
\centering
\includegraphics[width=1\textwidth]{images/Image4.png}
\caption{Fiche filière — Compétences et technologies du filière GTR}
\label{fig:fiche_comp}
\end{figure}
% ----------------------------------------------------------
\subsection{Connexion}
% ----------------------------------------------------------

\textbf{Route :} \texttt{/login}

\textbf{Overview :}
Page d'authentification réservée aux utilisateurs disposant d'une adresse académique valide,
permettant l'accès sécurisé à l'espace personnel ou au panel d'administration.

\textbf{Key Elements :}
\begin{itemize}
    \item Champ de saisie de l'adresse e-mail académique (\texttt{@uca.ac.ma} ou \texttt{@ensas.uca.ma}).
    \item Champ de saisie du code d'accès (haché via \textit{scrypt} côté serveur).
    \item Message d'erreur contextuel en cas d'identifiants invalides ou de compte désactivé.
    \item Redirection automatique vers \texttt{/dashboard} (étudiant) ou \texttt{/admin/analytics} (admin).
\end{itemize}

\textbf{Interface Description :}
La page présente un formulaire centré avec deux champs de saisie et un bouton de soumission.
La validation du domaine e-mail est effectuée côté serveur avant toute interrogation de la base.
Les erreurs sont affichées inline sous le formulaire via le mécanisme \texttt{flash} de Flask.
Le design est épuré, avec un fond neutre et un logo de l'établissement en en-tête. La figure~\ref{fig:login} illustre la page de connexion.

\begin{figure}[H]
\centering
\includegraphics[width=\textwidth]{images/Image5.png}
\caption{Page de connexion — formulaire d'authentification académique}
\label{fig:login}
\end{figure}

% ----------------------------------------------------------
\subsection{Test d'orientation}
% ----------------------------------------------------------
\subsubsection{Test d'orientation — Introduction}
\textbf{Route :} \texttt{/test-orientation}

\textbf{Overview :}
Page statique affichée avant le démarrage du test, fournissant à l'étudiant les informations
nécessaires sur le déroulement du questionnaire et l'invitant à commencer via un appel à l'action.

\textbf{Key Elements :}
\begin{itemize}
    \item Présentation du test et de son objectif (évaluation de la compatibilité filière).
    \item Instructions générales sur le format des réponses (échelle de Likert 1 à 5).
    \item Informations sur le déroulement : nombre de questions, durée estimée, confidentialité.
    \item Bouton de démarrage déclenchant le chargement du questionnaire interactif.
\end{itemize}

\textbf{Interface Description :}
La page adopte une mise en page centrée et épurée, sans aucun élément de questionnaire affiché.
Elle se compose d'un bloc d'information structuré (titre, description, liste des consignes) et
d'un bouton de démarrage bien visible, positionné en bas de page. L'objectif de cette interface
est de préparer l'étudiant à la passation, en assurant clarté, lisibilité et réduction de l'anxiété
cognitive avant d'entrer dans le questionnaire. La figure~\ref{fig:test_intro} montre la page d'introduction avec les consignes détaillées.

\begin{figure}[H]
\centering
\includegraphics[width=\textwidth]{images/Image6.png}

\vspace{0.5cm}

\includegraphics[width=\textwidth]{images/Image7.png}
\caption{Page d'introduction au test d'orientation — consignes et lancement du questionnaire}
\label{fig:test_intro}
\end{figure}

% ----------------------------------------------------------
\subsubsection{Test d'orientation — Questionnaire}
% ----------------------------------------------------------

\textbf{Route :} \texttt{/test-orientation}

\textbf{Overview :}
Interface interactive de passation du questionnaire psychométrique, permettant à l'étudiant
de répondre à chaque question et d'obtenir ses résultats de compatibilité par filière à l'issue du test.

\textbf{Key Elements :}
\begin{itemize}
    \item Questionnaire paginé : une question affichée par vue, avec navigation séquentielle.
    \item Mélange aléatoire reproductible des questions (seed fixe par session).
    \item Échelle de réponse de type Likert (1 à 5) sous forme de boutons sélectionnables.
    \item Barre de progression dynamique indiquant l'avancement dans le questionnaire.
    \item Affichage des résultats : classement des six filières par score de compatibilité (\%) et profil calculé.
\end{itemize}

\textbf{Interface Description :}
Pendant la passation, chaque question est rendue individuellement avec ses cinq boutons de réponse.
La barre de progression est mise à jour dynamiquement via JavaScript à chaque question validée.
À la soumission finale, les réponses sont transmises en JSON à la route \texttt{/api/submit-test} ;
les résultats sont restitués sur la même page sous forme d'un classement ordonné des six filières
avec leur pourcentage de compatibilité et le type de profil déterminé (\textit{Spécialiste},
\textit{Hybride} ou \textit{Équilibré}). La figure~\ref{fig:test_quest1} et la figure~\ref{fig:test_quest2} présentent l'interface du questionnaire et la barre de progression.

\begin{figure}[H]
\centering
\includegraphics[width=\textwidth]{images/Image8.png}
\caption{Interface du questionnaire du test d'orientation — navigation par question et barre de progression}
\label{fig:test_quest1}
\end{figure}

\begin{figure}[H]
\centering
\includegraphics[width=\textwidth]{images/Image9.png}
\caption{Interface du questionnaire du test d'orientation — navigation par question et barre de progression}
\label{fig:test_quest2}
\end{figure}

% ----------------------------------------------------------
\subsection{Tableau de bord}
% ----------------------------------------------------------

\textbf{Route :} \texttt{/dashboard}

\textbf{Overview :}
Espace personnel de l'étudiant authentifié, centralisant l'historique de ses tests et
présentant un profil global consolidé calculé sur l'ensemble de ses résultats.

\textbf{Key Elements :}
\begin{itemize}
    \item Liste chronologique des tests passés avec date, durée et résultats associés.
    \item Profil global agrégé (moyenne filière par filière sur tous les tests).
    \item Type de profil consolidé : \textit{Spécialiste}, \textit{Hybride} ou \textit{Équilibré}.
    \item Bouton d'accès rapide pour repasser le test d'orientation.
\end{itemize}

\textbf{Interface Description :}
Page sécurisée côté serveur, organisée en deux zones : un panneau supérieur affichant le profil global avec les scores moyens par filière, et un panneau inférieur présentant l’historique des tests avec la date, la durée, le profil obtenu et des détails extensibles. Le profil global est recalculé dynamiquement à chaque chargement. La figure~\ref{fig:dashboard_ui} illustre le tableau de bord avec l'historique et le profil global consolidé.
\begin{figure}[H]
\centering
\includegraphics[width=1\textwidth]{images/Image10.png}
\caption{Tableau de bord étudiant — historique des tests et profil global consolidé}
\label{fig:dashboard_ui}
\end{figure}

% ----------------------------------------------------------
\subsection{Administration — Analytiques}
% ----------------------------------------------------------

\textbf{Route :} \texttt{/admin/analytics}

\textbf{Overview :}
Interface de supervision statistique réservée à l'administrateur, offrant une vue globale
sur les orientations produites par la plateforme et la distribution des profils étudiants.

\textbf{Key Elements :}
\begin{itemize}
    \item Histogramme de répartition des recommandations par filière.
    \item Diagramme en anneau illustrant la distribution des profils (Spécialiste, Hybride, Équilibré).
    \item Courbe de tendance mensuelle du volume de tests réalisés.
    \item Filtre temporel dynamique via paramètre \texttt{?days=N}.
    \item Export CSV des données brutes via \texttt{/admin/export\_analytics\_csv}.
\end{itemize}

\textbf{Interface Description :}
La page est structurée en un tableau de bord à plusieurs panneaux graphiques rendus via
des bibliothèques JavaScript de visualisation. Un sélecteur de période permet de restreindre
l'analyse à une fenêtre glissante. Les trois graphiques sont disposés en grille et mis à jour
dynamiquement selon le filtre sélectionné. L'accès à cette page est conditionné par le
décorateur \texttt{@admin\_required}, qui redirige tout utilisateur non habilité. La figure~\ref{fig:admin_analytics} affiche le tableau de bord analytique de l'administration.

\begin{figure}[H]
\centering
\includegraphics[width=\textwidth]{images/Image11.png}
\caption{Administration — tableau de bord analytique avec indicateurs statistiques globaux}
\label{fig:admin_analytics}
\end{figure}

% ----------------------------------------------------------
\subsection{Administration — Étudiants}
% ----------------------------------------------------------

\textbf{Route :} \texttt{/admin/users}

\textbf{Overview :}
Interface de gestion des comptes étudiants, permettant à l'administrateur d'effectuer
l'ensemble des opérations CRUD sur les utilisateurs de la plateforme.

\textbf{Key Elements :}
\begin{itemize}
    \item Liste paginée des comptes étudiants avec initiales, filière conseillée et nombre de tests.
    \item Formulaire d'ajout de compte (nom complet, e-mail, code d'accès initial).
    \item Actions par compte : activation/désactivation, réinitialisation du code, suppression.
    \item Indicateur visuel du statut du compte (actif / désactivé) et de la filière recommandée.
\end{itemize}

\textbf{Interface Description :}
La page présente un tableau récapitulatif des comptes enregistrés, chaque ligne exposant les
informations essentielles de l'étudiant et un groupe d'actions rapides. Les opérations de gestion
sont réalisées de manière asynchrone via l'API REST dédiée (\texttt{POST}, \texttt{DELETE}).
Une garde logicielle bloque toute tentative de suppression d'un compte administrateur.
Le formulaire d'ajout est accessible via un panneau latéral ou une modale, sans rechargement
de page. La figure~\ref{fig:admin_users} présente la gestion des comptes étudiants, tandis que la figure~\ref{fig:admin_add} montre le formulaire d'ajout d'un compte.

\begin{figure}[H]
\centering
\includegraphics[width=\textwidth]{images/Image12.png}
\caption{Administration — gestion des comptes étudiants (ajout, activation, suppression)}
\label{fig:admin_users}
\end{figure}

\begin{figure}[H]
\centering
\includegraphics[width=\textwidth]{images/Image13.png}
\caption{Administration — Ajout d'un compte étudiant}
\label{fig:admin_add}
\end{figure}
% ----------------------------------------------------------
\subsection{Chatbot (composant global)}
% ----------------------------------------------------------

\textbf{Route :} Composant flottant — visible sur toutes les pages.

\textbf{Overview :}
Widget conversationnel intégré à l'ensemble de l'interface, permettant à l'utilisateur d'obtenir
des réponses contextuelles d'orientation sans quitter la page en cours.

\textbf{Key Elements :}
\begin{itemize}
    \item Bouton flottant (icône de bulle de dialogue) ancré en bas à droite de chaque page.
    \item Fenêtre de chat expansible avec historique de la conversation en cours.
    \item Champ de saisie de message et bouton d'envoi.
    \item Réponses générées par le modèle \texttt{mistral-tiny} via l'API Mistral AI.
\end{itemize}

\textbf{Interface Description :}
Le widget est implémenté comme un composant HTML/CSS/JavaScript indépendant, inclus dans le
template de base \texttt{base.html} et donc hérité par toutes les pages. Il s'affiche en
superposition (\texttt{position: fixed}) sans perturber la mise en page sous-jacente. La
communication avec le serveur s'effectue de manière asynchrone via \texttt{fetch} vers la route
\texttt{/api/chat}. Les messages sont affichés dans une fenêtre scrollable alternant bulles
utilisateur et bulles assistant. En cas d'erreur API (timeout, HTTP), un message de repli
explicite est affiché sans interrompre l'expérience. La figure~\ref{fig:chatbot} expose le widget de chat accessible sur toutes les pages de l'application.

\begin{figure}[H]
\centering
\begin{minipage}{0.48\textwidth}
    \centering
    \includegraphics[width=\textwidth]{images/Image14.png}
\end{minipage}
\hfill
\begin{minipage}{0.48\textwidth}
    \centering
    \includegraphics[width=\textwidth]{images/Image15.png}
\end{minipage}
\caption{Chatbot — widget flottant d’assistance à l’orientation, accessible sur toutes les pages}
\label{fig:chatbot}
\end{figure}
\newpage
% ============================================================
\clearpage
\addcontentsline{toc}{section}{Conclusion générale et perspectives}
\section*{Conclusion générale et perspectives}
% ============================================================

\textbf{1. Objectifs du projet} \\
Le projet \textbf{ENSA Guide} a été conçu avec pour objectif principal d'accompagner les étudiants du cycle préparatoire de l'ENSA Safi dans le choix de leur filière d'ingénierie. Face à l'absence d'un outil centralisé, l'application propose une évaluation objective des profils étudiants via un questionnaire psychométrique basé sur les aptitudes comportementales, afin de réduire l'incertitude liée à ce choix crucial.

\vspace{0.5cm}
\textbf{2. Contexte et outils utilisés} \\
Réalisé dans le cadre de notre formation, ce projet a mobilisé une architecture MVC robuste. Le développement s'est appuyé sur Python et le framework Flask pour la logique métier, couplé à SQLite pour la gestion de la base de données. L'interface utilisateur, développée en HTML, CSS et JavaScript vanilla, garantit une navigation intuitive. De plus, l'intégration de l'API Mistral AI a permis de doter la plateforme d'un assistant virtuel intelligent, offrant un accompagnement personnalisé en langage naturel.

\vspace{0.5cm}
\textbf{3. Limites du travail} \\
Bien que fonctionnelle, la plateforme présente certaines limites techniques et ergonomiques. Le chargement des feuilles de style (CSS) dépend parfois de chemins relatifs, ce qui peut causer des problèmes d'affichage sur certaines routes. De plus, la synchronisation de la base de données nécessite actuellement une réinitialisation manuelle lors de modifications des données initiales. Enfin, l'interface n'est pas encore totalement adaptative (responsive) sur l'ensemble des résolutions mobiles.

\vspace{0.5cm}
\textbf{4. Améliorations possibles (Perspectives)} \\
Plusieurs perspectives d'évolution sont envisageables pour pallier ces faiblesses et enrichir l'application. Techniquement, l'implémentation de Flask-Migrate permettrait une gestion fluide des versions de la base de données. Sur le plan ergonomique, une refonte \textit{mobile-first} des interfaces garantirait une accessibilité universelle. Fonctionnellement, il serait pertinent d'enrichir le contexte de l'intelligence artificielle avec des données plus exhaustives sur les filières, et d'ajouter un module d'exportation PDF pour les bilans d'orientation individuels.

\vspace{0.5cm}
\textbf{5. Bilan personnel} \\
Sur le plan personnel, ce projet s'est avéré extrêmement formateur. Il nous a permis de consolider nos compétences en développement web \textit{full-stack}, de l'architecture logicielle jusqu'à l'intégration d'intelligences artificielles génératives. Au-delà de l'aspect technique, le travail en équipe et la gestion rigoureuse d'un planning sous forme de diagramme de Gantt ont renforcé notre capacité à mener un projet d'ingénierie de bout en bout, tout en répondant à un besoin réel de notre environnement académique.

% ============================================================
\clearpage
\addcontentsline{toc}{section}{Bibliographie et Nétographie}
\section*{Bibliographie et Nétographie}
% ============================================================

\begin{thebibliography}{99}
    \bibitem{flask} \textbf{Documentation Flask} : \textit{Flask (A Python Microframework)}, Pallets Projects. Disponible sur : \url{https://flask.palletsprojects.com/} — Consulté le : 15 Mai 2026.
    \bibitem{sqlite} \textbf{Documentation SQLite} : \textit{SQLite Documentation}. Disponible sur : \url{https://www.sqlite.org/docs.html} — Consulté le : 15 Mai 2026.
    \bibitem{python} \textbf{Python Official Documentation} : \textit{Python 3 Documentation}. Disponible sur : \url{https://docs.python.org/3/} — Consulté le : 15 Mai 2026.
    \bibitem{mistral} \textbf{Mistral AI} : \textit{Documentation officielle de l'API Mistral AI}. Disponible sur : \url{https://docs.mistral.ai/} — Consulté le : 15 Mai 2026.
    \bibitem{jinja} \textbf{Jinja2} : \textit{Jinja2 Templating Documentation}. Disponible sur : \url{https://jinja.palletsprojects.com/} — Consulté le : 15 Mai 2026.
    \bibitem{numpy} \textbf{NumPy} : \textit{NumPy Documentation}. Disponible sur : \url{https://numpy.org/doc/} — Consulté le : 15 Mai 2026.
    \bibitem{werkzeug} \textbf{Werkzeug} : \textit{Werkzeug Documentation}. Disponible sur : \url{https://werkzeug.palletsprojects.com/} — Consulté le : 15 Mai 2026.
    \bibitem{mdn} \textbf{MDN Web Docs} : \textit{HTML, CSS, JavaScript}, Mozilla. Disponible sur : \url{https://developer.mozilla.org/} — Consulté le : 15 Mai 2026.
    \bibitem{mvc} \textbf{Architecture MVC} : \textit{Modèle-vue-contrôleur}, notions et principes de conception logicielle.
    \bibitem{ensa} \textbf{ENSA Safi} : Livret de l'étudiant et ressources internes de l'École Nationale des Sciences Appliquées de Safi (descriptions des filières et programmes).
\end{thebibliography}

\end{document}